\documentclass[11pt]{article}
\usepackage[margin=0.8in]{geometry}
\usepackage[T1]{fontenc}
\usepackage{lmodern}
\usepackage{amsmath,amssymb,amsthm,booktabs,array}
\usepackage[colorlinks=true,linkcolor=blue,citecolor=blue,urlcolor=blue]{hyperref}
\setlength{\parindent}{0pt}
\setlength{\parskip}{0.38em}
\setlength{\emergencystretch}{2em}
\Urlmuskip=0mu plus 2mu
\newtheorem{theorem}{Theorem}[section]
\newtheorem{lemma}[theorem]{Lemma}
\newtheorem{corollary}[theorem]{Corollary}
\newcommand{\N}{\mathbb{N}}
\newcommand{\Z}{\mathbb{Z}}
\newcommand{\Q}{\mathbb{Q}}
\newcommand{\R}{\mathbb{R}}

\title{Batch 16: Euler polynomial numerators and repeated Catalan convolutions}
\author{BlakeStone\\\small Proofs developed with ChatGPT}
\date{November 28, 2026}
\begin{document}
\maketitle
\begin{abstract}
We prove that $(2n-1)^2$ divides the reduced numerator of
$E_{2n-1}(n)$ for every positive integer $n$, the assertion recorded
in A291897. The proof is a specialization of an odd-degree centered
Euler-polynomial divisibility lemma. We also derive Mathar's A104722
recurrence from the literal self-convolution of repeated Catalan numbers,
with the necessary starting index $n\ge4$ and an explicit failure at $n=3$.
These are two mathematical families with four checked Lean modules.
Classical Euler and Catalan formulas are credited. The analytic
identification of the finite Euler formula is explained mathematically;
a separate analytic equivalence is outside the Lean development.
\end{abstract}

\section{A291897: squares in Euler-polynomial numerators}

Write $\operatorname{num}(q)$ for the numerator of a rational number $q$
when its denominator is positive and numerator and denominator are coprime.
The source entry A291897 is
\[
 a(n)=\operatorname{num}\bigl(E_{2n-1}(n)\bigr),\qquad n\ge1,
\]
where $E_m(z)$ denotes the Euler polynomial. Its recorded conjecture is
that $(2n-1)^2$ divides $a(n)$ \cite{a291897-source}.

Let $e_m$ be the signed Euler numbers. We use the conventions
\[
 \operatorname{sech}t=\sum_{m\ge0}e_m\frac{t^m}{m!},\qquad
 \frac{2e^{zt}}{e^t+1}=\sum_{m\ge0}E_m(z)\frac{t^m}{m!}.
\]
Thus $e_0=1$, $e_m=0$ for odd $m$, and for positive even $m$,
\[
 e_m=-\sum_{j=0}^{m-1}\binom mj e_j.
\]
The recurrence follows by comparing coefficients in
$\cosh t\operatorname{sech}t=1$ and proves that every $e_m$ is an integer.
These conventions and formulas are classical \cite{a291897-dlmf,a291897-carlitz}.

Define the integral centered polynomial
\begin{equation}\label{eq:euler-centered}
 P_m(x)=\sum_{j=0}^m\binom mj e_{m-j}x^j.
\end{equation}
Multiplying the exponential generating series gives
\[
 \sum_{m\ge0}P_m(x)\frac{t^m}{m!}
 =e^{xt}\operatorname{sech}t
 =\frac{2e^{(x+1)t}}{e^{2t}+1}
 =\sum_{m\ge0}2^m E_m\!\left(\frac{x+1}{2}\right)\frac{t^m}{m!}.
\]
Consequently
\begin{equation}\label{eq:euler-convention}
 P_m(x)=2^mE_m\!\left(\frac{x+1}{2}\right).
\end{equation}
This is also the centered finite expansion in DLMF~24.2.10.

\begin{theorem}\label{thm:euler-general}
Let $m\ge1$ be odd, and let $x\in\Z$ satisfy $m\mid x$. Then
\[
 m^2\mid P_m(x),\qquad
 m^2\mid\operatorname{num}\!\left(\frac{P_m(x)}{2^m}\right).
\]
\end{theorem}

\begin{proof}
In \eqref{eq:euler-centered}, the term $j=0$ is $e_m=0$.
The term $j=1$ is $m e_{m-1}x$, which is divisible by $m^2$.
For $j\ge2$, the factor $x^j$ is divisible by $m^2$.
Thus every summand is divisible by $m^2$, proving the first assertion.

Let $P_m(x)/2^m=u/v$ in lowest terms, with $v>0$. Cancellation
divides the original numerator and denominator by a common positive
integer $c$ dividing $2^m$. Since $m$ is odd, $\gcd(m^2,c)=1$.
The relation $m^2\mid P_m(x)=cu$ therefore implies $m^2\mid u$.
\end{proof}

\begin{corollary}\label{thm:euler-source}
For every integer $n\ge1$,
\[
 (2n-1)^2\mid\operatorname{num}\bigl(E_{2n-1}(n)\bigr).
\]
\end{corollary}

\begin{proof}
Set $m=x=2n-1$ in Theorem~\ref{thm:euler-general} and use
\eqref{eq:euler-convention}.
\end{proof}

\paragraph{Formal scope.}
\texttt{EulerSource.lean} defines the signed Euler numbers by the integer
recurrence, then uses the finite centered expansion to define the rational
Euler value and its reduced numerator. \texttt{EulerProof.lean} proves
the square-divisibility theorem, the coprime cancellation lemma, and
Corollary~\ref{thm:euler-source} for every positive natural index.
All divisibility and rational-reduction steps are formalized. The identification
with the analytic generating-function convention is the classical calculation
above; a separate analytic equivalence is not a Lean theorem in this package.

\paragraph{Attribution and selected scope.}
The existing admitted open-math-agents statement uses the same finite
source representation and credits \texttt{@DhruRules}; that representation
is acknowledged \cite{a291897-oma}. The classical centered expansion is
not a new formula. Shevelev's source-linked paper proves a different
denominator theorem for $E_n(x)-E_n(1)$ \cite{a291897-shevelev}.
The bounded prior-work screening is recorded in the accompanying package;
it does not establish worldwide priority. The general theorem and its
A291897 specialization are counted as one family. The adjacent
A292911 conjecture about a third factor of $2n-1$ is outside this result.


\section{A104722: a recurrence for repeated Catalan numbers}

Write $C_r=\frac1{r+1}\binom{2r}{r}$ for the Catalan numbers, and
define the repeated sequence and its self-convolution by
\begin{equation}\label{eq:a104722-source}
 b_n=C_{\lfloor n/2\rfloor},\qquad
 a_n=\sum_{k=0}^{n}b_kb_{n-k}\qquad(n\ge0).
\end{equation}
This is the literal definition of A104722 \cite{a104722-source}.
The initial values are $a_0,a_1,a_2,a_3=1,2,3,4$.

\begin{theorem}\label{thm:a104722}
For every integer $n\ge4$,
\[
 (n+4)a_n+(n+1)a_{n-1}
 -4(n+1)a_{n-2}+4(2-n)a_{n-3}=0.
\]
At $n=3$, the left side is $4$, so the starting index cannot be
lowered to $3$.
\end{theorem}

\begin{proof}
The standard Catalan convolution and ratio identities are
\begin{equation}\label{eq:a104722-catalan}
 \sum_{j=0}^{r}C_jC_{r-j}=C_{r+1},\qquad
 (r+2)C_{r+1}=2(2r+1)C_r.
\end{equation}
The first is the coefficient form of $C(x)=1+xC(x)^2$;
the second follows from the binomial formula.

For the odd argument $2r+1$, pair the even and odd summation
indices in \eqref{eq:a104722-source}. Both terms in the $j$th
pair are $C_jC_{r-j}$, giving
\begin{equation}\label{eq:a104722-odd}
 a_{2r+1}=2\sum_{j=0}^{r}C_jC_{r-j}=2C_{r+1}.
\end{equation}
For the even argument $2r+2$, the even terms sum to
$\sum_{j=0}^{r+1}C_jC_{r+1-j}=C_{r+2}$ and the odd terms sum to
$\sum_{j=0}^{r}C_jC_{r-j}=C_{r+1}$. Hence
\begin{equation}\label{eq:a104722-even}
 a_{2r+2}=C_{r+2}+C_{r+1}.
\end{equation}
Both identities hold for every $r\ge0$.

Let $R(n)$ denote the recurrence's left side and set
\[
 H_r=(r+2)C_{r+1}-2(2r+1)C_r.
\]
The second identity in \eqref{eq:a104722-catalan} says $H_r=0$.
Substitute \eqref{eq:a104722-odd} and \eqref{eq:a104722-even}
into $R(n)$ and collect the three adjacent Catalan numbers.
For the two parity cases this gives the exact certificates
\begin{align*}
 R(2r+4)&=2H_{r+2}+6H_{r+1},\\
 R(2r+5)&=6H_{r+2}+2H_{r+1}.
\end{align*}
For clarity, their expanded forms are
\[
\begin{aligned}
R(2r+4)
 &=(2r+8)C_{r+3}+(-2r-2)C_{r+2}-(24r+36)C_{r+1},\\
R(2r+5)
 &=(6r+24)C_{r+3}-(22r+54)C_{r+2}-(8r+12)C_{r+1}.
\end{aligned}
\]
Thus both parity cases vanish and exhaust all integers $n\ge4$.
At the preceding index, direct substitution gives
\[
 R(3)=7\cdot4+4\cdot3-16\cdot2-4\cdot1=4.
\]
\end{proof}

\paragraph{Prior work and scope.}
The source records Mathar's conjecture on November 9, 2012, but
prints no lower bound on $n$ \cite{a104722-source}; the theorem
supplies the necessary bound. This is a single recurrence family.
The Catalan expressions are established background. The OEIS
entry credits James Sellers's May 19, 2008 generating-function
argument for agreement with A059348 after the first term.
The related sequence has a different OEIS offset
\cite{a104722-related}. Zhao, Li and Sun's Theorem 3.7
\cite{a104722-zhao} gives a related permutation enumeration in
Catalan numbers; no novelty is claimed here for those formulas.

The contribution of this package is the explicit recurrence
derivation from the literal convolution, its endpoint correction,
and its checked Lean proof. The bounded prior-work search is
documented separately and does not establish global priority.
The formalization uses mathlib's natural-number Catalan sequence,
proves both parity formulas by finite-sum pairing, derives the ratio
from the central-binomial-coefficient identity, and checks both
linear-combination certificates. Its final results are
\texttt{A104722.source\_recurrence} and
\texttt{A104722.index\_three\_exception}.

\section{Verification and scope of the release}

The accompanying repository packages contain all Lean sources, pinned
toolchain and dependency manifests, frozen OEIS source snapshots with
hashes, supplementary exact-arithmetic scripts, build logs, and separate
agent review reports. Each package's complete set of local proof modules
is listed as an explicit default build target. The command
\texttt{lake build} checks the package under Lean~4.19.0 and mathlib
\texttt{c44e0c8ee63ca166450922a373c7409c5d26b00b}.

For this release, each of the two packages passed an author build and
a separate fresh source-only build: four successful package builds in
total. The builds used separate local directories and copied no compiled
proof objects from one another. They shared the same pinned compiler and
external dependency caches. The word ``independent'' describes this
separate compilation and a review by another proof agent, not verification
by a second theorem prover or external peer review.

The printed principal theorem audits use only Lean's standard
\texttt{propext}, \texttt{Classical.choice}, and \texttt{Quot.sound}
axioms. Source checks exclude local admissions, custom axiom declarations,
and \texttt{native\_decide}. Input hashes in the build logs bind the
checks to the published sources. The supplementary finite calculations
help detect source, index, and sign mistakes; the universal mathematical
assertions are established by the human and Lean proofs, with the
formal boundaries stated in the relevant sections.

The counts in this collection refer to mathematical families and selected
recorded assertions. Supporting identities, boundary corrections, and
corollaries are not separately counted as solved open problems. Each
package's \texttt{source\_checks.md} records a bounded prior-work inquiry
and credits known ingredients. Neither the recorded conjecture label nor
those searches establishes current worldwide open status or priority.

\begin{thebibliography}{99}
\raggedright
\bibitem{a291897-source}
V.~Shevelev, OEIS A291897, \emph{Numerator of $E(2n-1,n)$}, 2017.
Frozen OEIS data at commit \texttt{091a90f51479669655f4e46343867ea2abd5d4ff};
\url{https://oeis.org/A291897}.

\bibitem{a291897-dlmf}
NIST Digital Library of Mathematical Functions,
formulas 24.2.6--24.2.10 and 24.5.4,
\url{https://dlmf.nist.gov/24.2} and \url{https://dlmf.nist.gov/24.5}.

\bibitem{a291897-carlitz}
L.~Carlitz, \emph{A Note on Euler Numbers and Polynomials},
Nagoya Mathematical Journal \textbf{7} (1954),
\url{https://doi.org/10.1017/S0027763000018043}.

\bibitem{a291897-shevelev}
V.~Shevelev, \emph{On a Luschny question}, arXiv:1708.08096v3 (2017),
\url{https://arxiv.org/abs/1708.08096v3}.

\bibitem{a291897-oma}
GuillaumeBAUD-coder/open-math-agents,
\texttt{statements/oeis-a291897.lean}, admitted statement crediting
\texttt{@DhruRules}, commit
\texttt{5494b7a26d84a2fe7ab23ca2e45e1e864f5714f9},
\url{https://github.com/GuillaumeBAUD-coder/open-math-agents}.


\bibitem{a104722-source}
P.~Barry; recurrence conjecture by R.~J.~Mathar,
\emph{OEIS A104722: Self-convolution of repeated Catalan numbers}.
The entry includes James Sellers's May 19, 2008 argument.
Immutable source commit
\texttt{091a90f51479669655f4e46343867ea2abd5d4ff}.
\url{https://oeis.org/A104722}.

\bibitem{a104722-related}
The OEIS Foundation, \emph{A059348}.
Immutable source commit
\texttt{091a90f51479669655f4e46343867ea2abd5d4ff}.
\url{https://oeis.org/A059348}.

\bibitem{a104722-zhao}
T.~Zhao, X.~Li and Y.~Sun,
\emph{Enumeration of Permutations Without Consecutive Descents
Avoiding Patterns of Length 3},
Advances in Mathematics (China) \textbf{54}(1) (2025), 50--60,
Theorem 3.7. \url{https://doi.org/10.11845/sxjz.2023026b}.

\end{thebibliography}
\end{document}
